\section{The complete small endomorphism spectrum}\label{sec:spectrum}

The spectral argument uses the exact equivariance identity
\eqref{eq:exact-action}, separately from the moment estimate.
We first keep the reference metric, allowing nonintegrable ambient
coefficients.

\begin{proposition}\label{prop:reference-spectrum}
For every sufficiently small $x\in V_{\det}$, the self-adjoint
zero-form operator $D_a^*D_a$ in the reference metric has exactly $N$
eigenvalues below one fixed threshold $\gamma>0$ and
\begin{equation}\label{eq:reference-spectrum}
 (1-Cr^2)\nu_j(x)\le\lambda_j^0(x)\le
               (1+Cr^2)\nu_j(x),\qquad 1\le j\le N.
\end{equation}
Its kernel is exactly $\ker C_x$ inside $P_0$, and the remaining
spectrum is bounded below by $\gamma$. At an integrable point this
is the endomorphism Dolbeault spectrum.
\end{proposition}

\begin{proof}
Put $A=[a,\cdot]$ and $D=d+A$. The decomposition of zero-forms is
$P_0\oplus Q_0L^2$. On $P_0$,
\[
 D\xi=-J_xC_x\xi,\qquad d^*D\xi=0.
\]
Let
\[
 G_f=P_0D^*DP_0,\qquad B_f=Q_0D^*DP_0,\qquad
 L_Q=Q_0D^*DQ_0 .
\]
Harmonic orthogonality and \eqref{eq:canonical-estimates} imply
\begin{equation}\label{eq:relative-block}
 q_x\le G_f\le(1+Cr^2)q_x,\qquad
 B_f=Q_0A^*DP_0,\qquad
 \|B_f\xi\|^2\le Cr^2\langle G_f\xi,\xi\rangle .
\end{equation}
The last inequality is relative to the energy of the same vector.
An absolute bound by $Cr^4\|\xi\|^2$ would be insufficient.

By \eqref{eq:complement-coercive}, $L_Q$ has a fixed positive
lower bound $\gamma_0$. Its form domain is $Q_0W^{1,2}$ and its
operator domain is $Q_0W^{2,2}$. Its resolvent is uniformly bounded
from $L^2$ to $W^{2,2}$ for $0\le\lambda\le\gamma_0/2$.
The finite Schur form is
\[
 S_\lambda=G_f-\lambda I
              -B_f^*(L_Q-\lambda)^{-1}B_f,
\]
and \eqref{eq:relative-block} gives
\begin{equation}\label{eq:schur}
 (1-Cr^2)G_f-\lambda I\le S_\lambda\le G_f-\lambda I .
\end{equation}
Completion of the square on $P_0\oplus Q_0W^{1,2}$ identifies the
negative index of $D^*D-\lambda$ with that of $S_\lambda$.
The triangular change is a bijection on this form domain, since the
resolvent correction lies in $Q_0W^{2,2}$.

As $\|G_f\|=O(r^2)$, the form $S_{\gamma_0/2}$ is negative
definite for small $r$. Hence exactly $N$ eigenvalues lie below
$\gamma=\gamma_0/2$. The two form inequalities in \eqref{eq:schur},
negative-index counting and finite min--max give
\eqref{eq:reference-spectrum}. This reasoning counts multiplicities
and includes zeros; it never divides by $\nu_j$.

If $C_x\xi=0$, exact equivariance gives $D\xi=0$.
Conversely a kernel element has parallel part $\xi$ with $q_x(\xi)=0$
by \eqref{eq:schur} at $\lambda=0$. Then $B_f\xi=0$, and positivity
of the complementary operator forces its complementary part to vanish.
The kernel is therefore literally $\ker C_x\subset P_0$.
\end{proof}

This is the relative Schur mechanism used for repeated-factor
extensions in \cite[Section~5]{wuebben4} and the full local model
of \cite[Section~5]{wuebben7}. Here canonical harmonic equivariance
supplies its hypotheses on a general background.

\subsection{Transfer to the actual HE metric}
Let $y=gx$ be the correction in \eqref{eq:S-bound}, with $g=e^S$.
Conjugation $T_g\xi=g\xi g^{-1}$ preserves the complex rank-weighted
trace-free algebra and satisfies
\[
 C_yT_g\xi=g(C_x\xi)g^{-1}.
\]
Both energy and mass are therefore bounded in both directions by
factors $e^{\pm4\|S\|}$, and generalized min--max gives
\begin{equation}\label{eq:finite-transfer}
 e^{-8\|S\|}\nu_j(x)\le\nu_j(y)\le e^{8\|S\|}\nu_j(x).
\end{equation}
The spectral algebra is complex. There is no requirement that
nonunitary conjugation preserve its Hermitian real subspace.

Changing from $h_0$ to $H_{t,y}$ changes squared endomorphism norms
by $1+O(\delta)$. Changing $\omega_0$ to $\omega_t$ changes both
zero-form masses and one-form energies by $1+O(|t|)$.
Their common form domain is $W^{1,2}$, and the differential is unchanged.
The Rayleigh quotient and min--max therefore transfer all ordered
eigenvalues relatively. The scalar determinant correction has no
effect on endomorphism norms. Finally, holomorphic pullback of both
bundle and metric in \eqref{eq:actual-metric} is an exact isometry
intertwining the Laplacians. Proposition~\ref{prop:reference-spectrum}
and \eqref{eq:finite-transfer} prove
\eqref{eq:spectrum-main}, with a common complementary gap after
decreasing the cutoff. The kernel identification is transported by
the same maps.

For every induced endomorphism metric the splitting
\[
 \End E_x=\C I\oplus\End_0E_x
\]
is orthogonal and preserved by the differential and its adjoint.
On the scalar part $D_a(fI)=(\ddbar f)I$.
The scalar Dolbeault Laplacian on the compact connected
$(X,\omega_t)$ has a one-dimensional constant kernel and a
uniformly positive first nonzero eigenvalue for small $t$.
Taking the smaller of the scalar and trace-free cutoffs completes
the full endomorphism statement of Theorem~\ref{thm:main}.

